\section{معماری و انتقال به میکروی دیگر}
فایل‌های تولیدشونده در
\lr{Core}
و لایهٔ برد در
\lr{Application}
قرار دارند. پروتکل عمومی و فرمان‌های سنسور در
\lr{Libraries}
هستند. فایل‌های سازنده در
\lr{Drivers}
برای هر دو ابزار مشترک‌اند.
\cite{cubef1}
 فایل شروع اسمبلی کامپایلرهای مختلف جداست و
ابزارهای همگام‌سازی و آزمون در پوشه‌های مستقل قرار دارند.

\subsection{قرارداد لایهٔ سخت‌افزار}
ساختار
\lr{OneWire\_Bus}
یک اشاره‌گر زمینه و callbackهای پایین‌کشیدن، رهاکردن، خواندن و تأخیر
میکروثانیه‌ای دریافت می‌کند. تأخیر باید با وقفهٔ بسته هم اجرا شود. توابع ورود
و خروج ناحیهٔ حساس اختیاری ولی جفت هستند؛ ورود وضعیت قبلی وقفه را برمی‌گرداند
و خروج دقیقاً همان وضعیت را بازیابی می‌کند. حذف هر دو فقط وقتی مجاز است که
محیط برنامه زمان‌بندی را به روش دیگری تضمین کند.

در نمونهٔ فعلی شمارندهٔ سیکل
\lr{DWT}
پس از تنظیم کلاک فعال می‌شود. نوشتن رجیسترهای GPIO نیاز به تغییر مکرر حالت
پایه ندارد. اختلاف شمارنده برای تأخیرهای کوتاه با سرریز سازگار است. میکروهای
فاقد شمارندهٔ قابل‌استفاده باید تأخیر را با تایمر پیاده کنند؛ هدرهای کتابخانه
هیچ وابستگی به STM32 ندارند.

\subsection{حافظه و هم‌زمانی}
برنامه ساختار مدیریت، دو آرایهٔ جداگانهٔ آدرس، آرایهٔ دما و وضعیت را با ظرفیت
یکسان تخصیص می‌دهد. همهٔ این بافرها باید تا پایان کار زنده بمانند. تخصیص پویا
و نمونهٔ سراسری اجباری وجود ندارد. دسترسی به یک باس باید سری شود؛ از وقفه
فراخوانی نکنید و هنگام تبادل کلاک را عوض نکنید. در سیستم عامل، قفل دسترسی باید
کل تراکنش را پوشش دهد و فقط محدود به یک callback نباشد.
